ここでは、個別の高被引用論文が長期的にどのように影響を持つかを分析するため、オブソレッセンスに注目し寿命構造を観察する。まず、オブソレッセンス状態への到達を判定するために、P.D.B Paroloら\cite{Par1}の長期被引用モデルおよびD. Wangら\cite{Wan1}の被引用減衰モデルを前提とし、次の3つの条件を設定する： 1. 出版年から十分な期間がある、2. ピーク期が十分に過ぎ去っている、3. 被引用が十分に減衰している。これら3つの条件の具体的な設定値を図\ref{fig:gr6}に示す。
